% =============================================================================
% Appendix: dataset revision (simplified version).
%
% Condensed version of appendix_dataset_revision.tex; labels are distinct, so
% both files can be included in the same document. Usage: \input after
% \appendix. Required packages: booktabs. Assumes UTF-8 input and T1 font
% encoding (both provided by the ACL template). No figures, no bibliography.
%
% Every number comes from config/removed_questions.json, output/04 - revised,
% output/06 - revised, the abstention subset and the exports of the annotation
% runs under output/knowledge-annotation-work/, and is printed by
% ablations/revision/make_summary.py, which also checks the artifacts against
% each other. The knowledge-area appendix is mentioned in prose only and never
% referenced with \ref.
% =============================================================================

\section{Dataset Revision}
\label{app:rev-short}

Nineteen questions (0.046\% of the 41,653 questions of the filtered corpus)
were removed because the task they pose cannot be recovered from the dataset:
the statement was replaced or truncated during extraction or pre-processing,
the supporting text, chart or table it refers to was lost, or the item belongs
to a block that can only be solved jointly. As evaluation items they could
only be answered by chance. The defects were found as a by-product of the
knowledge-area annotation, whose annotator may abstain, and completed by a
string-pattern sweep. Nothing was removed for difficulty, for an unexpected
label or because the annotation passes disagreed.

\paragraph{Summary.} (1)~Five defect families, 19 questions from six exams
(Table~\ref{tab:rev-short-removed}). (2)~Ten were abstentions of the
annotation run with taxonomy~1.2, seven were found by a sweep for the same
patterns, two were abstentions of the final run. (3)~Look-alike groups were
read and kept: 78 of 80 questions with an ``INSTRUÇÕES'' header, 19 whose
last choice carries the heading of the next passage, and two whose support is
missing but whose task is identifiable. (4)~Removal is applied by a content
hash, never by row index, into a new directory with a manifest, a removal log
and a reload check, the original stages are never modified. (5)~The filtered
corpus goes from 41,653 to 41,634 questions and the deduplicated corpus from
37,887 to 37,873 (Table~\ref{tab:rev-short-datasets}).

\paragraph{Defects and criterion.} The second phase of the 2012 Informatics
Olympiad opens with three items whose statement is only ``A resposta da
questão $N$ é:'' (``The answer to question $N$ is:'') with the letters A to E
as choices, they form one puzzle and are undecidable as independent rows, in
two booklets (six rows). In seven BLUEX items the statement is a
pre-processing instruction (``Seu único objetivo é detectar as expressões
matemáticas, convertê-las para o padrão LaTeX\ldots'') with a worked example,
and the choices of the original literature, history or grammar question are
orphaned. Two POSCOMP 2024 items kept only an ``INSTRUÇÕES'' header and the
final question, without the data that define the computation time. Item~36 of
IME 2010 and 2011 asks ``What task below could Lammert B. Otten be legally in
charge of?'' without its reading passage, the two rows are identical and carry
different answer keys. COMVEST 2017 item~63 and ENEM 2018 item~136 refer to a
chart or table that is missing. A question is removed on two grounds only:
the statement or its supporting material is not recoverable from the dataset,
or the item depends on a block the one-row-per-question format does not
preserve.

\begin{table}[t]
\centering
\scriptsize
\setlength{\tabcolsep}{3pt}
\resizebox{\linewidth}{!}{
\begin{tabular}{@{}lp{3.0cm}rrl@{}}
\toprule
Family & Items & S04 & S06 & Detection \\
\midrule
Self-referential block & OBI 2012, phase 2, levels 1 and 2, items 1--3 & 6 & 3 & abst.\ 1.2 \\
Statement replaced & BLUEX UNICAMP 2018 (2), 2020 (68), 2023 (9), 2024 (7, 12); USP 2020 (37), 2022 (65) & 7 & 6 & abst.\ 1.2 (1), sweep (6) \\
Truncated statement & POSCOMP 2024 (19, 20) & 2 & 2 & abst.\ 1.2 (1), sweep (1) \\
Missing text & IME 2010 and 2011 (36) & 2 & 1 & abst.\ 1.2, run 1.4 \\
Missing data & COMVEST 2017 (63); ENEM 2018 (136) & 2 & 2 & abst.\ 1.4 \\
\midrule
Total & & 19 & 14 & \\
\bottomrule
\end{tabular}
}
\caption{The removed questions, with item numbers in parentheses. \emph{S04}, \emph{S06}: rows removed from the filtered (41,653) and the semantically deduplicated (37,887) corpora; five listed questions were already absent from the latter because the exact deduplication had removed them as copies of another listed question. \emph{Detection}: abstention of the full annotation run with taxonomy 1.2 or 1.4, or the pattern sweep.}
\label{tab:rev-short-removed}
\end{table}

\paragraph{Identification.} The full annotation run with taxonomy~1.2 ended
with 23 abstentions on 41,653 questions. The experimental taxonomy~1.3 run on
them labelled 14 and left 9, of which 8 had no recoverable statement and one
(AFA 2019 item~45, an English item without its passage but with an
identifiable task) was kept. The 8 were used as string patterns over the
filtered and deduplicated stages: the instruction string matches exactly 7
rows, 6 of which run~1.2 had labelled from their orphaned choices, ``A
resposta da questão'' matches exactly the 6 OBI rows, and the ``INSTRUÇÕES''
header matches 80 rows, of which only POSCOMP item~19 shares the truncation
of item~20. The list reached 15 questions on 7~October~2026. Taxonomy~1.4,
run on the 15 remaining abstentions, labelled the 2010 copy of the IME item
and abstained on the identical 2011 copy, the oscillation confirmed the
dependence on the missing passage and both were removed (17). The filtered
corpus revised with those 17 (41,636 rows) was the base of the final full
run, which ended with 8 abstentions, none repeated from earlier runs: 6
received a subject by hand and 2, whose chart or table is missing, were
removed on 8~October~2026 (19). The revised corpora were then regenerated
from the original stages with the complete list.

\paragraph{Validation.} Groups sharing a surface feature with a removed
question were read and kept: 80 rows starting with ``INSTRUÇÕES'' or
``Instruções:'' (65 POSCOMP 2024, 13 BACEN, one each ENADE 2006 and 2007),
all of which carry the full statement after the header except the two
POSCOMP items, 19 rows of COMVEST 2011--2015 and BLUEX USP 2020 and 2024
whose last choice carries, after its real text, the heading of the next
passage, which remain answerable and are recorded as a formatting clean-up,
and AFA 2019 item~45 and COMVEST 2015 item~38, where the competence assessed
is recognizable despite the missing support.

\paragraph{Procedure and effect.} Each question is identified by the SHA-256
of the canonical JSON of its exam, edition, item number, statement and
choices, the identifier the annotation uses, so one list serves every stage
and subset. The revision tool reads a dataset saved to disk, refuses an
existing output directory, never modifies its input, writes the revised
dataset with a removal log (input row index, identifier, exam, edition, item,
category, reason) and a manifest (hashes of input, output and list, removed
and absent identifiers), aborts if the metadata of a row differs from its
list entry, and reloads the output to verify that the kept rows are identical
to the input rows in the same order. Publishing re-hashes the output and
records the Hub revision in the manifest. The deduplicated corpus loses 14
questions: the other five were already removed by the exact deduplication as
copies of another listed question (the level-2 OBI items, UNICAMP 2024
item~12 and the 2010 IME copy). The annotated source corpus has the same
identifiers, in the same order, as the revised filtered corpus.

\begin{table}[t]
\centering
\scriptsize
\setlength{\tabcolsep}{3pt}
\resizebox{\linewidth}{!}{
\begin{tabular}{@{}lrrrl@{}}
\toprule
Dataset & Input & Rem. & Output & Revision \\
\midrule
Filtered, 17-question list & 41,653 & 17 & 41,636 & \textcolor{red}{\texttt{0e8b69a1}}\ldots\ (superseded) \\
Filtered, 19-question list & 41,653 & 19 & 41,634 & \textcolor{red}{\texttt{e0f638e0}}\ldots \\
Deduplicated & 37,887 & 14 & 37,873 & \textcolor{red}{\texttt{54ba11ff}}\ldots \\
Abstentions of run 1.2 & 23 & 8 & 15 & -- \\
Annotated export of run 1.4 & 41,636 & 2 & 41,634 & \textcolor{red}{\texttt{fb9d509a}}\ldots \\
\bottomrule
\end{tabular}
}
\caption{Datasets revised with the list. \emph{Rem.}: rows removed. Revisions are those of \textcolor{red}{\texttt{bench-temp-2/mmlu-pt-revised} (filtered), \texttt{bench-temp-2/mmlu-pt-deduplicated} and \texttt{bench-temp-2/mmlu-pt-knowledge-annotated}, pushed on 8~October~2026}; the abstention subset is a local working dataset.}
\label{tab:rev-short-datasets}
\end{table}

\paragraph{Published splits.} The final test questions are selected from
the 37,873-row revised deduplicated corpus. Their annotations come from the
41,634-row annotated source corpus: each selected question is joined by
\texttt{annotation\_id}, with unique identifiers, complete annotation
coverage and equality of all original question fields. The academic level
and macro area define the configurations. Each configuration has five dev
examples with answer rationales; their identifiers are excluded globally
from the test set. The published splits contain 37,798 test questions and 
75 dev examples (Table~\ref{tab:rev-short-splits}), with preserved source
annotations and deterministic record order. Evaluation uses these test
splits in the zero-shot and five-shot protocols, with the dev examples
providing the demonstrations for the latter.

\begin{table}[t]
\centering
\footnotesize
\begin{tabular}{@{}lrrl@{}}
\toprule
Level & Test & Dev & Revision \\
\midrule
High school & 8,704 & 30 & \texttt{39387109}\ldots \\
Undergraduate & 29,094 & 45 & \texttt{026191be}\ldots \\
\midrule
Total & 37,798 & 75 & -- \\
\bottomrule
\end{tabular}
\caption{Published splits of \texttt{bench-temp-2/mmlu-pt-high-school} and
\texttt{bench-temp-2/mmlu-pt-undergraduate}, respectively. Dev supplies
few-shot demonstrations; test contains the evaluation questions.
The full revisions are
\texttt{3938710959f204ae54ed}\allowbreak\texttt{257fcfe4ffed7c459046} and
\texttt{026191bef1e9690c0dc1}\allowbreak\texttt{18c81ef6ab427d285bf6}.}
\label{tab:rev-short-splits}
\end{table}

\paragraph{Limitations.} The detector is the abstention of one annotation
model plus string patterns derived from its abstentions, so only defects that
make the model abstain or share a pattern with one of them were found. The
corpus was not audited question by question. The engineering audit of the
first full run tagged 23 questions as source issues and 8 remain: six with a
missing photograph, chart or text but an identifiable task, and two not
re-examined under the removal criterion (BLUEX UNICAMP 2023 item~3, whose
statement field holds an explanatory text rather than a question, and ENEM
2020 item~1, whose poster is missing). ``Unrecoverable'' was judged against
the dataset, not the source PDFs. Identical copies were caught only because
their labels diverged.
